% -*- coding: utf-8 -*-
% !TEX program = xelatex
\documentclass[plain,noanswer,medmath]{../../template/jnuexam}
\usepackage{../../template/kysx}

\begin{document}


\examtitle{1990}{5}


\loadproblems{4}{1990P4}%载入试卷四


\makepart{填空题}{本题满分15分，每小题3分}


% 使用试卷四，第一大题，第1小题
\useproblem{4}{1}{1}

\useproblem{4}{1}{2}

\useproblem{4}{1}{3}

\useproblem{4}{1}{4}


\begin{problem}
已知随机变量$X\sim N(-3,1)$，$Y\sim N(2,1)$，且$X$, $Y$相互独立，
设随机变量$Z=X-2Y+7$，则$Z\sim$ \fillin{}.
\end{problem}


\makepart{选择题}{本题满分15分，每小题3分}


\useproblem{4}{2}{1}

\useproblem{4}{2}{2}

\useproblem{4}{2}{3}


\begin{problem}
设$A$为$n$阶可逆矩阵，$A^*$是$A$的伴随矩阵，则$\left|A^*\right|=$\pickout{}
\begin{abcd}
\item $|A|^{n-1}$．
\item $|A|$．
\item $|A|^n$．
\item $|A|^{-1}$．
\end{abcd}
\end{problem}


\begin{problem}
已知随机变量$X$服从二项分布，且$EX=2.4$，$DX=1.44$，
则二项分布的参数$n$，$p$的值为\pickout{}
\begin{abcd}
\item $n=4$，$p=0.6$．
\item $n=6$，$p=0.4$．
\item $n=8$，$p=0.3$．
\item $n=24$，$p=0.1$．
\end{abcd}
\end{problem}


\makepart{计算题}{本题满分20分，每小题5分}


\begin{problem}
求极限$\lim_{x\to \infty}\,\frac{1}{x}\int_0^x (1+t^2)\e^{t^2-x^2} \dt$．
\end{problem}


\begin{problem}
求不定积分$\int \frac{x\cos^4\frac{x}{2}}{\sin^3 x}\dx$．
\end{problem}


\begin{problem}
设$x^2+z^2=y\phi \left(\frac{z}{y} \right)$，其中$\phi$为可微函数，求$\frac{\pdz}{\pdy}$．
\end{problem}


\useproblem{4}{3}{2}


\makepart{}{本题满分9分}


\useproblem{4}{4}{0}%【同试卷四第四题】


\makepart{}{本题满分6分}


\begin{problem}
证明不等式$1+x\ln(x+\sqrt{1+x^2}) \ge \sqrt{1+x^2}$（$-\infty <x<+\infty$）．
\end{problem}


\makepart{}{本题满分4分}


\begin{problem}
设$A$为$10\times 10$矩阵$\begin{pmatrix}
        0 &      1 &      0 &      0 & 0 \\
        0 &      0 &      1 &      0 & 0 \\
   \cdots & \cdots & \cdots & \cdots & \cdots \\
        0 &      0 &      0 &      0 & 1 \\
  10^{10} &      0 &      0 &      0 & 0
\end{pmatrix}$，计算行列式$|A-\lambda E|$，其中$E$为$10$阶单位矩阵，$\lambda$为常数．
\end{problem}


\makepart{}{本题满分5分}


\begin{problem}
设方阵$A$满足$A^TA=E$，其中$A^T$是$A$的转置矩阵，$E$为单位阵．
试证明$A$所对应的特征值的绝对值等于$1$．
\end{problem}


\makepart{}{本题满分8分}

\useproblem{4}{6}{0}%【同试卷四第六题】


\makepart{}{本题满分5分}

\useproblem{4}{9}{0}%【同试卷四第九题】


\makepart{}{本题满分6分}


\begin{problem}
甲乙两人独立地各进行两次射击，假设甲的命中率为$0.2$，乙的为$0.5$，
以$X$和$Y$分别表示甲和乙的命中次数，试求$X$和$Y$的联合概率分布．
\end{problem}


\makepart{}{本题满分7分}

\useproblem{4}{11}{0}%【同试卷四第十一题】


\end{document}
